\chapter*{Tóm Tắt}
\addcontentsline{toc}{chapter}{Tóm tắt}

\textbf{Bối cảnh:} Transcript do nhận dạng tiếng nói tự động tạo ra trong họp trực tuyến thường ngắn, nhiễu và có hiện tượng chuyển mã ngôn ngữ. Phân tích sắc thái cảm xúc trên nguồn dữ liệu này có thể hỗ trợ quan sát chất lượng trao đổi, nhưng không phải là cơ sở để suy luận hay chẩn đoán sức khỏe tâm thần của một cá nhân.

\textbf{Mục tiêu:} Khóa luận xây dựng và kiểm chứng một ứng dụng MVP xử lý transcript tiếng Việt trong họp trực tuyến theo hướng có thể truy vết bằng chứng, bảo vệ quyền riêng tư và cập nhật gần thời gian thực. Ứng dụng được thiết kế để nhận một bộ phân loại sentiment có thể thay thế; đầu ra dự kiến gồm ba lớp \textit{positive}, \textit{neutral}, \textit{negative}; lớp \textit{unknown} được giữ lại khi dữ liệu hoặc bộ phân loại không đủ điều kiện.

\textbf{Phương pháp:} Hệ thống tiếp nhận transcript từ Vexa, lưu phiên bản segment trong backend sở hữu dữ liệu, áp dụng quality gate và language router, ghép các segment cùng người nói thành analysis unit, sau đó gọi adapter phân loại và tổng hợp theo cửa sổ thời gian. Kiến trúc dùng FastAPI, PostgreSQL, Valkey Stream, worker tách riêng và SSE. PhoBERT và custom model là lựa chọn tích hợp về mặt lý thuyết \cite{phobert2020}; runtime đã kiểm chứng của MVP là \texttt{baseline-rules-v2} dùng để kiểm tra luồng xử lý, không phải kết quả fine-tune hay đánh giá chất lượng mô hình.

\textbf{Kết quả:} Luận văn chuẩn bị biểu mẫu kiểm chứng pipeline và bảng so sánh nhiều model để điền sau quá trình build, huấn luyện lại và chạy thực nghiệm. Tại thời điểm biên soạn, không công bố giá trị định lượng về transcript, coverage, chất lượng sentiment, WER ASR, độ trễ hay năng lực sàng lọc trầm cảm. Dashboard được thiết kế để hiển thị KPI, timeline, evidence và cập nhật SSE khi có dữ liệu chạy thực tế.

\textbf{Kết luận:} MVP tạo nền tảng có thể kiểm toán để tích hợp bộ phân loại sentiment vào transcript họp. Huấn luyện lại, benchmark, calibration và đánh giá trải nghiệm là hướng tối ưu hóa tiếp theo, không phải đóng góp thực nghiệm của khóa luận này.

\vspace{0.5cm}
\textbf{Từ khóa:} \textit{Phân tích sentiment, PhoBERT, transcript tiếng Việt, họp trực tuyến, event-driven, evidence, quyền riêng tư.}

\vspace{1cm}
\hrule
\vspace{0.5cm}

\chapter*{Abstract}
\addcontentsline{toc}{chapter}{Abstract}

\textbf{Background:} Automatic-speech-recognition transcripts from online meetings are often short, noisy, and code-switched. Sentiment analysis can assist the observation of conversational content, but it must not be used to infer or diagnose an individual's mental health.

\textbf{Objective:} This thesis builds and verifies an evidence-traceable, privacy-aware, near-real-time application MVP for Vietnamese meeting transcripts. The application is designed to accept a replaceable sentiment classifier. The intended task comprises \textit{positive}, \textit{neutral}, and \textit{negative} classes; \textit{unknown} is retained when data or classifier conditions are insufficient.

\textbf{Methods:} The system ingests Vexa transcripts, versions segments in an application-owned backend, applies a quality gate and language router, constructs speaker-bounded analysis units, invokes a classifier adapter, and publishes windowed aggregates. FastAPI, PostgreSQL, Valkey Stream, a separate worker, and SSE form the implementation. PhoBERT and custom models are theoretical integration options \cite{phobert2020}; the verified MVP runtime is \texttt{baseline-rules-v2}, which exercises the processing flow rather than claiming a fine-tuned-model evaluation.

\textbf{Results:} The thesis provides pipeline-verification and multi-model-comparison templates to be populated after build, retraining, and experimental runs. At the time of writing, no quantitative values for transcript volume, coverage, sentiment quality, ASR WER, or latency are reported. The dashboard is designed to display KPIs, a timeline, evidence, and SSE updates when actual run data are available.

\textbf{Keywords:} \textit{Sentiment analysis, PhoBERT, Vietnamese transcript, online meeting, event-driven architecture, evidence, privacy.}

\cleardoublepage
